\documentclass[UTF8,fontset=none,11pt,a4paper]{ctexart}
\usepackage{ipara-handbook}
\handbooksetup{NET-07 拥塞控制}{408 · 计算机网络 NET-07}{Feedback · cwnd · AIMD · Reno}
\providecommand{\tightlist}{\setlength{\itemsep}{0pt}\setlength{\parskip}{0pt}}
\providecommand{\passthrough}[1]{#1}
\begin{document}
\handbooktitle{拥塞控制}{在不知道路径容量时闭环试探}{408 · NET-07}
\section{本册定位}\label{0-ux672cux518cux5b9aux4f4d}

本 Topic
回答：大量独立发送者共享有限链路与队列、又看不到全局负载时，端点怎样用反馈调节
offered load，避免网络进入高排队、高丢包、低有效吞吐的失控状态？

本册拥有 congestion
signal、\passthrough{\lstinline!cwnd!}、\passthrough{\lstinline!ssthresh!}、slow
start、congestion avoidance、AIMD、fast retransmit/recovery 的 408
经典模型，以及效率、公平、时延和稳定性的权衡。

本册使用 TCP sequence/ACK/connection
state，见\href{../06_传输层_端点_UDP与TCP状态机/README.md}{传输层}；可靠性与接收端流控分别不等同于
congestion control。

\section{独立阅读词典：拥塞是怎样被观察和控制的}

\begin{handbooktable}{L{2.3cm} L{4.3cm} Y}
\toprule
术语 & 本册中的可操作定义 & 不要混成什么 \\
\midrule
congestion / queue & 工作量超过某个服务环节能及时处理；queue 是暂存等待的 packet 集合 & 都不是 packet 自带的标签 \\
offered load $\lambda$ / service rate $C$ & $\lambda$ 是 sender 试图注入的速率；$C$ 是瓶颈每秒能处理的量 & $\lambda$ 不是有效吞吐量 \\
cwnd / ssthresh & cwnd 限制在途数据；ssthresh 是从快增转慢增的本地分界估计 & 都不是路径容量真值 \\
slow start & 用 ACK 在 RTT 粒度快速扩大探测窗口 & “slow”不是线性增长 \\
congestion avoidance / AIMD & 接近可用范围后慢增；AIMD 是加性增大、乘性减小 & 不是停止发送或唯一实现 \\
loss/timeout / duplicate ACK & 丢包、超时和重复 ACK 是不同强度的路径反馈证据 & 都不是拥塞原因的完整证明 \\
\shortstack[l]{ECN/AQM/RED\\/goodput} & ECN 标记反馈；AQM 提前标记/丢弃；goodput 是交给应用的有效速率 & AQM 不是端点 cwnd 更新动作 \\
\bottomrule
\end{handbooktable}

因此，看到 loss 或 timeout 时，先把它当作路径反馈，再判断 sender 应如何改变 cwnd。
丢包是“网络已承受压力”的证据，不是拥塞这个原因本身；重传还可能把 offered load
进一步推高。

\section{根本问题：总输入可以超过共享服务能力}\label{1-ux6839ux672cux95eeux9898ux603bux8f93ux5165ux53efux4ee5ux8d85ux8fc7ux5171ux4eabux670dux52a1ux80fdux529b}

设某瓶颈链路服务率为 \(C\)，多个流的总到达速率为 \(\lambda\)：

\begin{itemize}
\tightlist
\item
  \(\lambda<C\)：队列通常能被清空；
\item
  \(\lambda\approx C\)：小波动就会积累排队；
\item
  \(\lambda>C\) 持续存在：队列增长，最终丢包或时延失控。
\end{itemize}

朴素方案是每个 sender 都尽可能快地发送。单个 sender
看起来在提高自身吞吐，但所有 sender 的独立最优叠加会破坏共享资源：

\[
\boxed{
\text{More offered load}
\to \text{Queue growth}
\to \text{Delay/loss}
\to \text{Retransmission}
\to \text{Even more load}
}
\]

拥塞控制必须打断这个正反馈。

\section{拥塞是路径状态，不是单个 packet
的属性}\label{2-ux62e5ux585eux662fux8defux5f84ux72b6ux6001ux4e0dux662fux5355ux4e2a-packet-ux7684ux5c5eux6027}

发送端通常不知道：瓶颈在哪里、可用容量多少、同时有多少竞争流。它只能观察反馈：

\begin{itemize}
\tightlist
\item
  loss/timeout；
\item
  duplicate ACK pattern；
\item
  increasing RTT/queue delay；
\item
  ECN 等显式标记；
\item
  ACK 到达节奏。
\end{itemize}

然后更新本地发送约束：

\[
\boxed{
\text{Probe}
\to \text{Observe feedback}
\to \text{Infer path state}
\to \text{Adjust cwnd}
\to \text{Probe again}
}
\]

这是反馈控制，不是一次算出带宽后永久使用。

\section{\texorpdfstring{\texttt{cwnd}
的对象与单位}{cwnd 的对象与单位}}\label{3-cwnd-ux7684ux5bf9ux8c61ux4e0eux5355ux4f4d}

\passthrough{\lstinline!cwnd!} 是 sender 为该 connection 维护的
congestion window，限制网络中允许存在的未确认数据量。TCP 实际发送还受
receiver window：

\[
FlightSize\le \min(cwnd,rwnd).
\]

408 题目有时用 MSS 个数表示窗口，有时用 byte。计算前必须写单位；不能把
\(\text{cwnd} = 8\) 自动解释成 8 byte 或 8 segment。

若 sender 希望充分利用路径，理想在途量与路径 BDP
同量级；大幅超过时，多余数据主要进入 queue，而不是让传播管道变长。

\section{Slow
Start：在未知容量下快速探测}\label{4-slow-startux5728ux672aux77e5ux5bb9ux91cfux4e0bux5febux901fux63a2ux6d4b}

连接开始或严重拥塞后，若从 1 MSS 线性增加，找到高速路径容量太慢。Slow
Start 借 ACK clock 让 \passthrough{\lstinline!cwnd!} 在每个 RTT 约翻倍：

\begin{lstlisting}
cwnd = 1 MSS
one RTT of ACKs -> about 2 MSS
next RTT        -> about 4 MSS
next RTT        -> about 8 MSS
...
\end{lstlisting}

名称中的 ``slow'' 是相对一开始直接发送巨大 burst；它的增长在 RTT
维度上是指数级。

当 \passthrough{\lstinline!cwnd!} 到达
\passthrough{\lstinline!ssthresh!} 附近，切换到 congestion
avoidance。\passthrough{\lstinline!ssthresh!}
不是路径真实容量，而是根据历史拥塞更新的本地分界估计。

\section{Congestion
Avoidance：加性试探}\label{5-congestion-avoidanceux52a0ux6027ux8bd5ux63a2}

接近已知可用范围后，继续指数增长会频繁过冲。经典 congestion avoidance 让
\passthrough{\lstinline!cwnd!} 每 RTT 大约增加 1 MSS：

\[
\Delta cwnd\approx 1\ MSS/RTT.
\]

每个 ACK 的具体增量可按当前 \passthrough{\lstinline!cwnd!}
分摊，使一整个窗口被确认后累计约 1
MSS。核心不是背某个实现公式，而是识别增长从 multiplicative 变为
additive。

\section{Congestion event
后为什么要乘性下降}\label{6-congestion-event-ux540eux4e3aux4ec0ux4e48ux8981ux4e58ux6027ux4e0bux964d}

若检测到拥塞仍只减一个固定小量，多个 flow
可能长时间把总负载维持在瓶颈之上。经典 AIMD：

\[
\boxed{\text{Additive Increase} + \text{Multiplicative Decrease}}
\]

用缓慢线性增长持续探测剩余容量，用比例下降快速释放大量在途负载。AIMD
在理想同质条件下还具有向效率与公平状态收敛的几何直觉，但真实
RTT、MSS、应用模式和算法差异会影响公平性。

\section{Timeout 与 duplicate ACK
表示不同强度的证据}\label{7-timeout-ux4e0e-duplicate-ack-ux8868ux793aux4e0dux540cux5f3aux5ea6ux7684ux8bc1ux636e}

\subsection{Timeout：反馈闭环严重中断}\label{71-timeoutux53cdux9988ux95edux73afux4e25ux91cdux4e2dux65ad}

RTO 到期意味着对应数据在足够长时间内没有获得预期确认。经典 408/Reno
模型把它视为较严重拥塞：更新 \passthrough{\lstinline!ssthresh!}，将
\passthrough{\lstinline!cwnd!} 降到较小初值，重新 slow start。

\subsection{Three duplicate
ACKs：后续数据仍在通过}\label{72-three-duplicate-acksux540eux7eedux6570ux636eux4ecdux5728ux901aux8fc7}

连续 duplicate ACK 表明接收端反复看见同一个缺口，但缺口之后的一些
segments 已到达。这通常比 timeout
提供更细的信息：路径仍在交付数据，只是某个 segment 可能丢失。

经典模型由此产生：

\begin{itemize}
\tightlist
\item
  fast retransmit：不等 RTO，重传推定丢失 segment；
\item
  fast recovery：不把发送速率退回最初状态，而按相对温和方式降低并继续。
\end{itemize}

具体 \(\text{cwnd}/\text{ssthresh}\) 数值更新必须服从题目指定的
Tahoe/Reno 教材模型；现代 TCP 实现有多种 congestion-control
algorithms，不能把一张 408 状态图冒充所有工程实现。

\section{ACK
clock：反馈不仅说``收到''，还提供节奏}\label{8-ack-clockux53cdux9988ux4e0dux4ec5ux8bf4ux6536ux5230ux8fd8ux63d0ux4f9bux8282ux594f}

当 ACK 随数据穿过瓶颈返回，ACK 到达节奏近似反映路径的交付节奏。sender 按
ACK 释放新数据，可以避免无节制 burst。

但 ACK compression、delayed ACK、reverse-path congestion
等会扭曲时钟。ACK clock 是有用反馈机制，不是精确的全局容量测量仪。

\section{Queue
的双重身份}\label{9-queue-ux7684ux53ccux91cdux8eabux4efd}

Queue 不是纯粹坏事：短队列吸收
burst，使链路持续忙碌。问题是持久过量队列：

\begin{itemize}
\tightlist
\item
  latency 增大；
\item
  timeout 和交互性能恶化；
\item
  buffer overflow 导致 loss；
\item
  新流和短流被已有大队列拖累。
\end{itemize}

因此优化目标不是``绝不排队''或``永不丢包''，而是在
utilization、delay、loss、fairness 和 stability 之间取舍。

\section{Network-assisted 与 end-to-end
feedback}\label{10-network-assisted-ux4e0e-end-to-end-feedback}

拥塞控制先分“事前约束”与“事后反馈”。事前/open-loop 策略在观察本次拥塞前限制系统行为，例如接纳控制、资源预留、流量整形、调度与丢弃策略；反馈/closed-loop 策略检测 queue、delay、loss 或显式标记，再把信号送回能降低 offered load 的主体。前者降低进入危险区的概率，后者适应实际运行状态；二者可以同时存在。

再按反馈 Owner 分层：router/network layer 可调度队列、丢弃或标记 packet，也可在特定体系中反馈显式速率；transport endpoint 根据 ACK/loss/ECN/RTT 等调整发送窗口或速率。网络层“发现/标记拥塞”不自动完成端点减速，端点算法也不能直接观察每台 router 的完整队列真相。

\begin{itemize}
\tightlist
\item
  loss-based control：从丢包/ACK 推断拥塞，部署简单但信号较晚；
\item
  ECN：网络设备在丢包前标记拥塞，端点仍负责降低发送；
\item
  delay-based control：从 RTT/queue delay
  变化提前推断，但需要区分路径基线和噪声；
\item
  explicit rate
  等更强网络协助模型直接给出速率或资源信息，但需要基础设施支持。
\end{itemize}

408 核心以经典 TCP loss-based
模型为主；扩展用于解释反馈设计空间，不替代考试指定算法。

因此网络层拥塞与 TCP 拥塞控制不是两种互斥现象：前者描述共享 router/link 的资源状态与网络可提供的控制动作，后者是端点借助有限反馈调节 offered load 的一种闭环实现。

\section{408 状态推进：把题设算法写成转移函数}

先固定单位（byte 或 MSS）、初值、`ssthresh`、ACK 粒度与题设版本。然后对每个事件应用：

\begin{handbooktable}{L{2.6cm} Y Y Y}
\toprule
事件 & 进入前状态 & 经典教材动作 & 停止/下一分支\\
\midrule
new ACK，$cwnd<ssthresh$ & slow start & 每 ACK 增长；一 RTT 近似翻倍 & 达阈值转 avoidance\\
new ACK，$cwnd\ge ssthresh$ & avoidance & 一 RTT 近似增加 1 MSS & 继续探测\\
timeout & 任一发送阶段 & $ssthresh\leftarrow$ 题设比例；$cwnd\leftarrow$ 小初值 & 重传并回 slow start\\
three duplicate ACKs，Tahoe & 检测单缺口 & fast retransmit；窗口降到小初值 & slow start\\
three duplicate ACKs，Reno & 检测单缺口 & fast retransmit；进入/维持 fast recovery & new ACK 后回 avoidance\\
\bottomrule
\end{handbooktable}

具体加一、减半、初始窗口和 recovery 膨胀数值以题目定义为准。手册拥有的是“状态、事件、反馈、转移”模型，不把某一版本常数写成所有 TCP 的永恒事实。

\begin{examplebox}{轮次与逐 ACK 不能混算}
若题目说 slow start 每收到一个确认新数据的 ACK 增加 1 MSS，则一个窗口内的多个 ACK 累计造成下一 RTT 可发送量近似翻倍；不能把“每 ACK 加 1”误写成“每 ACK 翻倍”。题目若直接按 RTT 给窗口序列，则不要再重复逐 ACK 累加。
\end{examplebox}

\section{拥塞性能核验层：负载增加为何可能降低吞吐}

“资源需求超过可用资源”只是起点。评价拥塞至少要同时观察 offered load、throughput/goodput、queue length、delay/jitter、drop/retransmission。典型负载曲线分为：
\begin{handbooktable}{L{2.8cm} L{3.0cm} Y}
\toprule
区间 & 曲线关系 & 系统含义\\
\midrule
未饱和 & throughput 近似随 offered load 线性增长 & 新输入大多转化为有效输出\\
轻度拥塞 & 输出仍增长，但斜率下降 & queue/drop 已开始消耗新增负载\\
拥塞 & offered load 增加，throughput 反而下降 & 重传、排队与丢弃进入正反馈\\
死锁极限 & goodput 接近 0 & 大量资源用于无效等待/重复工作\\
理想控制上界 & 达容量后近似水平 & 超额输入在破坏内部状态前被抑制\\
\bottomrule
\end{handbooktable}

因此 packet loss 是拥塞的\textbf{征兆/反馈载体}，不是拥塞的根本原因；真正原因是持续资源需求与服务能力不匹配。发送量或链路 utilization 变大也不自动等于性能更好，必须把 retransmission/header 从 application goodput 中扣除。

扩大单个 buffer 不能消除 $\lambda>C$。它只把早期 drop 改成更长 queue：原 segment 尚在排队时 sender 已因 RTO 重传，原件与副本同时消耗链路，进一步增大 offered load。提高单台 processor 或一条 link 也可能只把瓶颈移到下一资源。正确问题不是“哪个部件不够大”，而是“端到端服务链中哪个约束先饱和，反馈能否在正反馈形成前降低输入”。

\section{TCP 算法核验层：逐 ACK 公式与证据强度}

slow start 的逐 ACK 教材表达不是无条件加一个完整 MSS。若一个 ACK 新确认了 $N$ bytes，常用上界为
\[
\Delta cwnd\le\min\{N,\mathrm{SMSS}\}.
\]
当一个窗口的数据都被新 ACK 覆盖时，累积增量约等于原窗口，于是 RTT 粒度近似翻倍。Delayed ACK、ACK loss、application-limited 与 rwnd bottleneck 都可能让实际轨迹偏离理想 $1,2,4,8$；做题采用轮次近似时，应把这些前提显式写出。

congestion avoidance 若要在一个 RTT 内累计约 $1\,\mathrm{SMSS}$，可把每个新 ACK 的增量分摊为量级
\[
\Delta cwnd\approx\frac{\mathrm{SMSS}^2}{cwnd}.
\]
这解释了为什么“每 RTT 加 1 MSS”与“每 ACK 也会更新”并不矛盾。它们是同一过程的不同时间粒度。

三个 duplicate ACK 的证据来自缺口之后仍有 segment 到达：每个失序到达促使 receiver 重复声明同一累计 ACK。它比 timeout 更早，也证明至少还有数据穿过路径；但它仍是启发式阈值，packet reordering 也可能制造 duplicate ACK。阈值取 3 是误判乱序与等待 RTO 之间的工程折中，不是数学必然。

\begin{handbooktable}{L{2.7cm} L{2.7cm} Y Y}
\toprule
事件 & 共同动作 & Tahoe 教材模型 & Reno 简化模型\\
\midrule
timeout & 先按题设把 $ssthresh$ 更新为事件前窗口的比例 & $cwnd$ 回小初值，slow start & 同 Tahoe\\
3 duplicate ACKs & fast retransmit 推定缺口 & $cwnd$ 回小初值，slow start & $cwnd\leftarrow ssthresh$，进入 avoidance/recovery\\
\bottomrule
\end{handbooktable}

Reno 的更细实现常在进入 fast recovery 时暂时令 $cwnd=ssthresh+3\,\mathrm{SMSS}$，把三个 duplicate ACK 解释为三个 segment 已离开网络；后续 duplicate ACK 可继续膨胀窗口，填补缺口的新 ACK 到达后退出 recovery。408 简化题常直接写 $cwnd=ssthresh$。两者只可按题设择一，不能把 recovery 临时膨胀与下一轮稳定窗口重复相加。

\section{计算题口径契约：先定义状态时点，再开始算}

补充笔记中的细节材料 提供了一个内部一致的做题 convention，但它必须作为\textbf{题解口径}而非协议标准使用：
\begin{enumerate}
\item “第 $n$ 轮末发生事件”表示第 $n$ 轮仍用事件前窗口发送，转移动作从第 $n+1$ 轮可见；“当 $cwnd=X$ 时发生事件”则直接在该状态处理。
\item 教材对 $cwnd=ssthresh$ 允许 slow start 或 avoidance；本册默认按 avoidance，题设另定则覆盖。
\item 若题目以整 MSS 做离散轮次，Source 采用 slow start 到门限时封顶、减半向下取整；真实 byte-based 算法和其他教材可能不采用同一离散化，必须先声明。
\item 事件处理先保存 $cwnd_{event}$，再算新 $ssthresh$，最后设置新 $cwnd$ 与阶段；不能先改窗口再拿新值减半。
\item $rwnd$ 只约束实际 send window；除非题设显式改算法，拥塞事件仍更新 $cwnd/ssthresh$，不拿 $\min(cwnd,rwnd)$ 代替事件前 $cwnd$。
\end{enumerate}

\begin{examplebox}{Reno 全轨迹：两种事件怎样分叉}
设 $cwnd_0=1$ MSS、$ssthresh_0=8$ MSS；第 6 轮末发生 3 duplicate ACK，第 13 轮末 timeout。采用上述 convention：
\[
\begin{aligned}
&1,2,4,8,9,10,11\mid 5,6,7,8,9,10,11\mid1,2,4,5,6,\\
&\text{第 6 轮末： }ssthresh=\lfloor11/2\rfloor=5,\quad cwnd=5,\\
&\text{第 13 轮末： }ssthresh=\lfloor11/2\rfloor=5,\quad cwnd=1.
\end{aligned}
\]
第一条竖线是 Reno fast recovery 的温和退让，第二条是 timeout 后回 slow start。若改为 Tahoe，第一事件也令 $cwnd=1$，之后同一轮次的事件前窗口会改变，后续 $ssthresh$ 不能照抄 Reno。
\end{examplebox}

含恒定 $rwnd=12$、而算法状态 $cwnd=20$ 时发生 timeout，实际窗口可只有 12，但新门限仍按题设用事件前 $cwnd=20$ 计算为 10，而不是用 12 得 6。若问累计发送量，只有在“每轮能填满 $\min(cwnd,rwnd)$ 且 ACK 正常返回”的前提下，才可直接把各轮窗口求和；平均 goodput 还要除以总 RTT 时间并扣除重传。

反推曲线可用五个局部形状：翻倍近似 slow start；每轮 $+1$ 近似 avoidance；降到小初值多为 timeout/Tahoe；减半后线性增长多为 Reno duplicate-ACK 分支；既不翻倍也不加一的上升，若恰等于 $ssthresh$，可能是题目 convention 下的门限封顶。最后用“事件后门限来自事件前窗口”“Reno 简化 recovery 后窗口等于新门限”交叉核验。

\section{网络反馈核验层：Tail Drop、AQM 与 RED}

拥塞闭环要区分 signal generation 与 endpoint response。router 决定何时 drop/mark、影响哪些 flows；TCP sender 决定如何改 $cwnd$。只写“路由器丢包”没有说明信号质量。

FIFO tail drop 在 queue 满后连续丢弃新到 packet，容易同时打击共享队列中的许多 TCP flows。这些 flows 近乎同时降低窗口、随后又同步增长，形成 global synchronization，使链路负载大幅振荡。扩大 queue 可能推迟这一刻，却增加 standing delay 和 timeout 风险。

AQM 在 queue 真正满之前主动 drop 或 ECN-mark 少量 packet，目标是让反馈\textbf{更早、更分散}。RED 是教学中的一个实例：
\begin{handbooktable}{L{3.1cm} Y Y}
\toprule
平均队列长度 $\bar q$ & 动作 & 设计目的\\
\midrule
$\bar q<q_{min}$ & 正常入队 & 吸收短 burst\\
$q_{min}\le\bar q<q_{max}$ & 以随 $\bar q$ 增长的概率 drop/mark & 随机分散反馈，避免 flows 同步\\
$\bar q\ge q_{max}$ & drop/mark 强制动作 & 阻止队列继续失控\\
\bottomrule
\end{handbooktable}

判据使用 average queue 而非瞬时 queue，相当于过滤短 burst；中间区随机选择而非确定性全丢，是为了不把提前门限重新变成另一种 tail drop。RED 的具体参数很难调，不能从这个经典例子推出它是唯一或当前通用 AQM；稳定 Owner 是“主动队列管理改善反馈时点与分散度”的机制。

open-loop 与 closed-loop 也不是互斥协议。Admission control、reservation、traffic shaping、scheduler 等在运行事件前限制风险；closed-loop 执行 monitor $\to$ convey signal $\to$ adjust。抑制分组本身会增加负载，捎带 ECN 标记开销小但需端点配合，探测则有固定开销与时延。反馈太快可能振荡，太慢则来不及阻止正反馈，控制时间常数本身就是系统设计变量。

\section{概念边界}\label{11-ux6982ux5ff5ux8fb9ux754c}

\begin{handbooktable}{L{2.0cm} C{0.45cm} L{2.4cm} Y Y}
\toprule
概念 A & $\neq$ & 概念 B & 真正区别与题目信号 & 混淆后果\\
\midrule
Congestion & $\neq$ & Receiver overflow & path resources vs receiver buffer & $cwnd/rwnd$ 更新写反\\
$cwnd$ & $\neq$ & $rwnd$ & sender inference vs receiver advertisement & 状态 Owner 错\\
$ssthresh$ & $\neq$ & Bottleneck capacity & 阶段阈值 vs 外部动态容量 & 把阈值当带宽\\
Timeout & $\neq$ & Certain congestion & 强推断，仍可能是误码/长时延 & 把观测当事实\\
Duplicate ACK & $\neq$ & New-data ACK & 重复缺口 vs 推进连续前缀 & ACK 计数错误\\
Flow rate & $\neq$ & Goodput & 可含首部/重传 vs 应用有效输出 & 发送更多即更好\\
Slow start & $\neq$ & Linear growth & 低起点、RTT 维度近似指数增长 & 轮次推演错误\\
Classic Reno & $\neq$ & All modern TCP & 408 算法实例 vs 多种现实实现 & 旧图当唯一事实\\
AQM & $\neq$ & Endpoint response & router 发信号 vs sender 调窗口 & 跨层动作合成一步\\
\bottomrule
\end{handbooktable}

\section{做题调用协议}\label{12-ux505aux9898ux8c03ux7528ux534fux8bae}

\begin{enumerate}
\def\labelenumi{\arabic{enumi}.}
\tightlist
\item
  标出
  \passthrough{\lstinline!cwnd!}、\passthrough{\lstinline!ssthresh!}、\passthrough{\lstinline!rwnd!}
  的单位和初值；
\item
  画 RTT 轮次或逐 ACK 事件，按题目粒度更新；
\item
  判断当前处于 slow start、avoidance 还是 recovery；
\item
  区分 timeout 与 duplicate ACK 触发的教材分支；
\item
  计算实际发送上限时取 \(\min(\text{cwnd}, \text{rwnd})\) 并扣除
  in-flight；
\item
  明确题目采用 Tahoe、Reno 或自定义规则；
\item
  用 BDP、queue 和 goodput 检查窗口变化是否有物理意义。
\end{enumerate}

\section{贯穿母例：为什么 1、2、4、8
后不能永远翻倍}\label{13-ux8d2fux7a7fux6bcdux4f8bux4e3aux4ec0ux4e48-1248-ux540eux4e0dux80fdux6c38ux8fdcux7ffbux500d}

假设瓶颈在途与队列可暂时容纳约 10 MSS。sender 从 1 MSS slow start：

\begin{lstlisting}
1 -> 2 -> 4 -> 8: path may absorb
8 -> 16: offered flight overshoots available capacity
-> queue grows / loss signal appears
-> sender reduces cwnd and records ssthresh
-> later probes additively near the observed boundary
\end{lstlisting}

Slow start 解决``完全不知道容量时怎样快速接近''；AIMD
解决``接近边界后怎样持续试探并在过冲时退让''。二者不是两套互不相关的背诵规则，而是探索阶段不同。

\section{高频 First
Divergence}\label{14-ux9ad8ux9891-first-divergence}

\begin{itemize}
\tightlist
\item
  把 \passthrough{\lstinline!cwnd!} 当 receiver 通告值：状态 owner 错；
\item
  每收到一个 ACK 就让 \passthrough{\lstinline!cwnd!} 翻倍：混淆逐 ACK
  增量与逐 RTT 效果；
\item
  到 \passthrough{\lstinline!ssthresh!} 后停止增长：忘记 avoidance
  仍在探测；
\item
  timeout 与三次重复 ACK 使用同一分支：丢失证据强度分层；
\item
  只算发送速率不算 goodput：重传负载被当成有效数据；
\item
  未读题设就套 Reno：没有确认教材模型。
\end{itemize}

\section{一页压缩与复原问题}\label{15-ux4e00ux9875ux538bux7f29ux4e0eux590dux539fux95eeux9898}

\[
\boxed{
\text{Probe capacity}
\to \text{Read ACK/loss/ECN/delay}
\to \text{Infer congestion}
\to \text{Adjust cwnd}
\to \text{Probe again}
}
\]

\begin{enumerate}
\def\labelenumi{\arabic{enumi}.}
\tightlist
\item
  为什么 congestion control 必须是闭环而不是固定速率？
\item
  Slow Start 为什么既``慢''又近似指数增长？
\item
  AIMD 的加性与乘性分别解决什么？
\item
  duplicate ACK 为什么通常比 timeout 提供更细的路径信息？
\item
  Queue 在什么范围内有益，何时开始伤害系统？
\end{enumerate}

\section{来源与校正说明}\label{16-ux6765ux6e90ux4e0eux6821ux6b63ux8bf4ux660e}

\begin{itemize}
\tightlist
\item
  早期归档没有独立拥塞控制正文；本轮已核验用户提供的 新提供的中文补充笔记，而非从术语表复制；
\item
  \passthrough{\lstinline!transport/tcp-congestion.md!} 已吸收负载曲线、拥塞推断、逐 ACK 增长、Tahoe/Reno、fast retransmit/recovery 与 AIMD 设计理由；
\item
  \passthrough{\lstinline!transport/tcp-congestion-detail.md!} 已吸收事件时点、阈值等号、离散封顶/取整 convention、逐轮/反推/累计量自检；这些口径被明确限制为题解契约；
\item
  \passthrough{\lstinline!network/network-layer-functions.md!} 中 congestion、open/closed loop、tail drop、global synchronization、AQM/RED 部分归本册；route/forwarding 内容仍由 NET04/NET05 拥有；
\item
  \passthrough{\lstinline!cwnd!}、\passthrough{\lstinline!rwnd!}
  共同限制在途数据，以及 slow start/congestion avoidance/fast
  retransmit/recovery 的经典边界依据
  \href{https://www.rfc-editor.org/rfc/rfc5681.html}{RFC 5681} 校正；
\item
  具体窗口初值和现代算法会演进，本册将工程现实与 408
  指定模型分层，不把版本相关常数写成永恒机制。
\end{itemize}
\end{document}
